\section{Contracting plaque jets and exact normal coordinates}
\label{nf:section}

We record graph and log-density derivatives of unstable plaques
at their crossings with a stable axis. Under $f$, these finite
state records contract, although the unstable plaques themselves
expand. We construct smooth normal coordinates at each finite level,
compatible with forgetting higher derivatives. Their polynomial
dynamics supplies the finite algebraic data for the hull construction
in Section~\ref{hull:section}; the stable-plaque construction uses $f^{-1}$.

We work over the ergodic transformation $f=\phi_T$ chosen in the
preliminaries.  All assertions involving Lyapunov decompositions concern
a common invariant conull set.  The conull set is refined countably many
times, in particular once for each finite jet order.  No regularity of
the normal coordinates as functions of their centers will be used.

\subsection{The raw graph and log-density towers}

Choose a transverse section at $v$ and symplectic coordinates $(x,y)$,
with $x,y\in\R^m$, in which the central stable and unstable plaques
are $y=0$ and $x=0$, respectively.  The two plaques are smooth
transverse Lagrangians, so the relative Lagrangian neighborhood
construction provides these coordinates.  One may first straighten
one plaque and then subtract the differential defining the other
plaque as a graph.  Choose the height coordinate so that height zero
agrees with both central strong plaques.  The resulting section
maps, height functions, and axis parametrizations have bounded norms
at every fixed smooth order on uniform smaller charts.  This follows
from the uniform leafwise bounds and the usual local symplectic
coordinate construction; none of these choices needs to vary
smoothly with $v$.

An unstable plaque crossing the stable axis at $(a,0)$ is written
$x=P(y)$, where $P(0)=a$.  Its conormal form on the section is
\begin{equation}\label{nf:graph-form}
 e^{c(x,y)+r(y)}
 \bigwedge_{i=1}^m\left(dx_i-\sum_{q=1}^m
                    \partial_qP_i(y)\,dy_q\right).
\end{equation}
Fixed nonzero units and orientation signs are suppressed in this
formula.  They will be retained as constant factors under a change
of chart.  We choose the smooth background $c$ so that $r(0)=0$ for
every actual crossing and $r=0$ along the central unstable plaque.
Here is why this choice is available without transverse smoothness.
Along the unstable axis the actual form $u$ is smooth.  Along the
stable axis its coefficient in \eqref{nf:graph-form} is the
reciprocal, up to the fixed wedge-pairing constant, of the coefficient
of $w$ in $\bigwedge dy_i$.  That coefficient is smooth along the
stable plaque.  The two logarithms agree at the origin after choosing
the constant units.  Extend them by
$c(x,y)=c(x,0)+c(0,y)-c(0,0)$.  For the stable construction the
background can then be chosen as $-c$.  These backgrounds have
uniform bounds at each fixed order.

For $j\geq1$ define the raw state space $\mathcal E_j(v)$ by coordinates
\begin{equation}\label{nf:raw-state}
 \left(a,p_1,\ldots,p_j,s_1,\ldots,s_j\right),
 \qquad p_l=D^lP(0),\quad s_l=D^lr(0).
\end{equation}
Set $\mathcal E_0(v)=\R^m$, with coordinate $a$; write
$\pi_{j,j-1}$ for the linear map forgetting the highest tensors.
These are unrestricted symmetric graph and scalar tensors.
In particular, we impose no Lagrangian equation on alternative
states.  Actual states give a continuous graph over $a$ at every
finite level.  Continuity of all of their finite jets follows from
uniform leafwise derivative bounds and interpolation on plaques.
The actual central state is zero at every level.

\begin{lemma}[Raw dynamics and its weights]\label{nf:raw-dynamics}
At every level $j$, the section map of $f$ induces a smooth
invertible germ
$G_{v,j}:\mathcal E_j(v)\to\mathcal E_j(fv)$ fixing zero.  These maps
commute with the projections $\pi_{j,j-1}$.  Their base component
depends only on $a$; their other components are polynomials in the
hidden coordinates $(p_l,s_l)$ with coefficients smooth in $a$.
For fixed $j$ their degrees are bounded independently of the length
of a composition along the central orbit.  The same statements hold
for changes of adapted charts over the same stable plaque.

The raw maps and their local inverses have uniformly bounded norms
at every fixed order on suitable uniform raw domains.  Let
$-\chi_1,\ldots,-\chi_m<0$ be the stable Lyapunov exponents of $f$,
with multiplicities, and put $\chi_*=\min_i\chi_i>0$.  The derivative
cocycle at zero has only negative exponents.  In the quotient
$\ker\pi_{j,j-1}$ they are sums of $j+1$ stable exponents on graph
tensors and sums of $j$ stable exponents on log-density tensors.
Thus their negative exponents define positive weights, each at
least $\chi_*$, and the new weights at level $j$ are at least
$j\chi_*$.
\end{lemma}

\begin{proof}
Write the section map as $(A(x,y),B(x,y))$ and its height shift as
$\rho(x,y)$.  It carries the stable axis to the stable axis, so
$B(x,0)=0$ and $D_xB(x,0)=0$.  Symplecticity gives
\begin{equation}\label{nf:axis-derivative}
 D_yB(a,0)=D_xA(a,0)^{-t}.
\end{equation}
Indeed the derivative on that axis is block upper triangular, and
preservation of the pairing between the two coordinate planes gives
this identity.  Put
\[
 Y(y)=B(P(y),y),\qquad P_*(Y(y))=A(P(y),y).
\]
Then $Y(0)=0$ and $DY(0)=D_yB(a,0)$, independently of $DP(0)$.
Consequently formal inversion of the $j$-jet of $Y$ uses only powers
of the base-dependent invertible matrix $D_yB(a,0)$ as denominators.
Successively differentiating the displayed equations proves that
the transformed graph jet through order $j$ is polynomial in the
hidden tensors, with smooth coefficients in $a$.

For the coefficient, let $M_*=DP_*(Y)$ and $x=P(y)$.  Pullback on
the $x$-directions yields the determinant
$\det(D_xA-M_*D_xB)$.  Our height convention says that the pullback
of the target section form is $e^{-hT+h\rho}$ times the source
section form.  Hence the transformed log record is given by
\begin{align}
 r_*(Y)={}&r(y)+c(x,y)-c_*(A,Y)
 \label{nf:density-transport}\\
 &-\log|\det(D_xA-M_*D_xB)|-hT+h\rho(x,y)
       +\text{constant}.\notag
\end{align}
At $y=0$ the determinant is $\det D_xA(a,0)$, independent of every
hidden tensor.  The right side has value zero for actual states by
actual transport and the definition of the backgrounds.  It
therefore has value zero for every alternative with that base $a$.

There is no missing $(j+1)$st graph derivative in
\eqref{nf:density-transport}.  In $j$ derivatives of
$M_*D_xB(P(y),y)$, the term with all $j$ derivatives falling on
$M_*$ is multiplied by $D_xB(a,0)=0$.  Every surviving term uses
at most $j-1$ derivatives of $M_*$, hence at most $j$ derivatives
of $P_*$.  All derivatives of the logarithm have only powers of
the nonzero base determinant as denominators.  This proves closure
at order $j$ and polynomial dependence for the density jet.

This computation is a universal finite-jet formula: it applies to
any smooth symplectic section map preserving the stable axis, with
its smooth height shift.  A composition of any length is again such
a map.  The maximal number of occurrences of the hidden tensors in
this universal formula depends only on $j$, not on the section
map.  This proves the asserted uniform degree bound for arbitrary
compositions.  Applying the formula to the inverse map proves
invertibility and the same finite-level assertions for the inverse.
For one time step all required section derivatives and inverse
axis determinants are uniformly controlled on smaller compact
charts.  The universal formulas then prove all the asserted raw
finite-order bounds.

It remains to identify the linear diagonal blocks.  At the central
state write $S=D_xA(0,0)$ and $U=D_yB(0,0)=S^{-t}$.
The base derivative is $S$.  Modulo lower graph and density orders,
the highest graph tensor transforms by
\[
 p_j\longmapsto S\,p_j(U^{-1}\,\cdot,\ldots,U^{-1}\,\cdot),
\]
and the highest scalar tensor transforms by
\[
 s_j\longmapsto s_j(U^{-1}\,\cdot,\ldots,U^{-1}\,\cdot).
\]
The latter may also receive a linear contribution from $p_j$.
Thus these quotient diagonal blocks are, respectively,
$S\otimes\operatorname{Sym}^jS$ and
$\operatorname{Sym}^jS$, with a triangular coupling.  Their exponents
are exactly the sums stated in the lemma.  The multiplicative
ergodic theorem applied to the invariant order filtration gives
no other exponents.
Since $j\geq1$, every new weight is at least $j\chi_*$.
\end{proof}

The stable graph tower uses $f^{-1}$, exchanges $x$ and $y$, and uses
$w$ with its factor $e^{hT}$ and the background $-c$.  The same proof
therefore gives another tower with a strictly positive minimum
weight.  Subsequent time indices count iterates of the indicated
map, which is $f$ for the unstable-plaque tower and $f^{-1}$ for its
stable counterpart.

\subsection{Estimates and measurable choices}

At each finite level, the raw dynamics is contracting. To construct
compatible normal coordinates, we need estimates for its linear blocks
and a way to solve the nonresonant coefficient equations.
We collect these facts here, together with the invariant-subspace
splitting later used for polynomial relations of actual states.

\begin{definition}\label{nf:tempered}
A measurable function $C\geq1$ on the regular set is \emph{tempered} if
\[
 \lim_{|n|\to\infty}\frac{\log C(f^n v)}{|n|}=0.
\]
A positive radius is tempered when its reciprocal, increased to at
least one, is tempered.  An estimate with arbitrarily small exponential
slack means that for every $\epsilon>0$ its prefactor can be chosen
tempered, while its exponential rate increases by at most $\epsilon$.
\end{definition}

If $C$ is tempered, the function
\[
 C_\epsilon(v)=\sup_{n\in\Z} C(f^n v)e^{-\epsilon|n|}
\]
is finite and measurable, majorizes $C$, and satisfies
$e^{-\epsilon}C_\epsilon(v)\leq C_\epsilon(fv)
\leq e^\epsilon C_\epsilon(v)$.  We refer to this operation, and to its
reciprocal for radii, as slow regularization.  The multiplicative
ergodic theorem, with its tempered norm comparison, gives the following
form of the estimates we use.  If a linear block $A$ has the single
exponent $\lambda$, then, in the original finite-dimensional norms,
\begin{equation}\label{nf:lyapunov-bounds}
 C_\epsilon(v)^{-1}e^{(\lambda-\epsilon)n}\|z\|
 \leq\|A_v^{(n)}z\|
 \leq C_\epsilon(v)e^{(\lambda+\epsilon)n}\|z\|,
 \qquad n\geq0.
\end{equation}
The corresponding estimates hold for inverses and with $v$ replaced
by any $f^qv$.  Tensor products add exponents, and duals negate them.

We will also use finite triangular arrays whose diagonal blocks have
one and the same exponent $\lambda$.  Suppose their diagonal blocks
satisfy \eqref{nf:lyapunov-bounds} and their one-step off-diagonal
coefficients are tempered.  Their products satisfy the same estimates
with additional arbitrarily small slack.  Indeed, in a triangular
array of length $r$, each matrix entry of a product is a sum over paths
with at most $r-1$ off-diagonal jumps.  There are at most a polynomial
in $n$ many choices of jump times.  Each diagonal interval contributes
its exponential factor; the interval lengths sum to $n$ minus at most
$r-1$.  Slowly regularized bounds at jump times contribute at most
$C(v)e^{r\delta n}$ for any preassigned $\delta>0$.  Absorb this factor
and the polynomial in $n$ by choosing $\delta$ smaller than the desired
slack.  Apply the identical argument to the triangular inverse to
obtain the lower estimate.  This proof does not assume logarithmic
integrability of the additional off-diagonal coefficients.

\begin{lemma}[Splitting invariant subspaces]\label{nf:blocks}
Suppose a measurable invertible linear cocycle has an invariant direct
sum $E=\bigoplus_{i=1}^r E_i$, with distinct exponents
$\lambda_1<\cdots<\lambda_r$, and each $E_i$ satisfies
\eqref{nf:lyapunov-bounds}.  Every measurable invariant subspace $S$
satisfies, almost everywhere,
\[
 S=\bigoplus_{i=1}^r(S\cap E_i).
\]
No temperedness assumption on $S$ or on a measurable frame for $S$ is
required.  The statement also applies to a cocycle whose base
transformation is $f^{-1}$.
\end{lemma}

\begin{proof}
Use the measurable direct-sum norm
$\|\sum_i z_i\|_\oplus=\max_i\|z_i\|$ on $E$.
Only measurability of this norm is needed in the argument below;
no tempered comparison with the original ambient norm is required.
Write $E=E_<\oplus E_r$ and put $K=S\cap E_<$ and
$I=\operatorname{pr}_{E_r}S$.  In quotient norms, $S/K$ is the graph of
a measurable linear map $L:I\to E_</K$.  Its transport obeys
\[
 \|L_{f^n v}\|
 \leq C_\epsilon(v)^2
 e^{(\lambda_{r-1}-\lambda_r+2\epsilon)n}\|L_v\|.
\]
Here the operator on the quotient has norm at most that on $E_<$,
and the inverse on $I$ has norm at most the inverse on $E_r$.
Choose $2\epsilon<\lambda_r-\lambda_{r-1}$.  Thus
$\|L_{f^n v}\|\to0$ almost everywhere.  For every $a>0$, invariance of
the probability measure and dominated convergence give
\[
 \mu\{\|L\|>a\}
 =\int 1_{\{\|L_{f^n v}\|>a\}}\,d\mu(v)\longrightarrow0.
\]
Consequently $L=0$ almost everywhere, so
$S=(S\cap E_<)\oplus(S\cap E_r)$.  Induction proves the assertion.
\end{proof}

\begin{lemma}[A nonresonant homological equation]\label{nf:homological}
Let $L$ be an invertible finite-dimensional measurable cocycle with a
single exponent $\gamma\ne0$ and estimates
\eqref{nf:lyapunov-bounds}.  For a tempered measurable inhomogeneous
term $e$, the recurrence
\begin{equation}\label{nf:recurrence}
 q_{n+1}=L_nq_n+e_{n+1}
\end{equation}
has a unique tempered solution on every regular two-sided orbit.
The solution is measurable.  This conclusion also holds for the
single-exponent triangular blocks described above.
\end{lemma}

\begin{proof}
For $\gamma<0$ the solution is
\[
 q_n=e_n+\sum_{j\geq1}
       (L_{n-1}\cdots L_{n-j})e_{n-j};
\]
for $\gamma>0$ it is
\[
 q_n=-\sum_{j\geq0}
       (L_n^{-1}\cdots L_{n+j}^{-1})e_{n+j+1}.
\]
Both expressions solve \eqref{nf:recurrence} by shifting the summation
index.  To verify convergence and temperedness, choose
$0<3\delta<|\gamma|$.  For example, in the first expression the $j$th
term has norm at most
\[
 C_\delta(v)E_\delta(v)
 e^{2\delta|n|}e^{(\gamma+3\delta)j},
\]
after slow regularization of the linear estimates and the input.
The geometric series converges.  Since $\delta$ is arbitrarily small,
its sum is tempered; the second expression has the same proof with
$-\gamma$ in place of $\gamma$.  Partial sums are measurable.  The
difference of two solutions is a homogeneous orbit.  A nonzero such
orbit grows exponentially in the future if $\gamma>0$ and in the
past if $\gamma<0$, contradicting temperedness.
\end{proof}

For later recurrence arguments, every finite collection of measurable
finite norms, inverse minors, and positive radii can be bounded on a
set of positive measure.  Birkhoff's theorem then supplies returns
with positive frequency.  Such a set and its bounds can depend on
the finite orders under discussion.  Leafwise essential suprema give
measurable bounds for continuous leafwise extensions: a supremum over
a compact plaque equals the essential supremum over its interior,
after using a slightly larger plaque.  The same observation applies
to every fixed derivative.  Absolute continuity and a countable
collection of plaque boxes allow us to impose agreement almost
everywhere on both whole strong leaves through a typical center.
Flow boxes give the corresponding weak-leaf assertion.  These facts
do not furnish a radius common to infinitely many orders.

\subsection{Exact resonance normalization}

A graded space in this subsection is a finite direct sum
$E=\bigoplus_{d>0}E_d$.  Give a coordinate on $E_d$ weight $d$.
A polynomial map is \emph{weight preserving} if its $E_d$ component
is a sum of monomials of weight $d$.  Since all weights are positive,
such a polynomial has degree at most $d_{\max}/d_{\min}$.  Its
nonlinear $E_d$ component uses only variables of weights strictly
less than $d$; consequently an invertible linear diagonal gives a
polynomial inverse by induction in weight.

\begin{lemma}[Exact compatible normal coordinates]
\label{nf:normalization}
For the raw towers in Lemma~\ref{nf:raw-dynamics}, at every finite
level there are measurable smooth germ coordinates
$\Psi_{v,j}$, tangent to the identity, and weight-preserving
polynomial diffeomorphisms $N_{v,j}$ such that
\begin{equation}\label{nf:conjugacy}
 \Psi_{fv,j}\circ G_{v,j}=N_{v,j}\circ\Psi_{v,j}.
\end{equation}
Their linear part is the raw derivative, with its Oseledets grading.
The coordinates and normal maps commute with the tower projections.
For every fixed state level and every fixed derivative order, their
norms and the corresponding inverse norms have measurable tempered
bounds on measurable tempered positive radii.  At a fixed level the
coordinate change is smooth on a forward admissible domain; the
estimates at different orders do not require one radius common to
the entire infinite tower.
\end{lemma}

\begin{proof}
Fix a finite level and suppress its index.  The derivative
$A_v=DG_v(0)$ has a measurable Oseledets splitting by weights $d$,
with exponent $-d$.  Oseledets projections commute with every tower
projection: the latter intertwines the derivative cocycles, and a
vector with one exponent can project only to that exponent or to
zero.  Use adapted norms, temperedly equivalent to the raw norms,
in which $\|A_v\|\leq e^{-\chi_*/2}$, decreasing the exponent
slack if necessary.

We first construct the formal conjugacy.  Suppose it has been
constructed through ordinary degree $b-1$.  The discrepancy in
degree $b$ is a homogeneous polynomial.  Its component with output
weight $d$ and input weights $d_1,\ldots,d_b$ lies in the cocycle
of homogeneous maps with linear transport
\[
 Q\longmapsto A_v Q(A_v^{-1}\,\cdot,\ldots,A_v^{-1}\,\cdot).
\]
This cocycle has the single exponent
$d_1+\cdots+d_b-d$.  If it is nonzero, Lemma~\ref{nf:homological}
solves the homological equation by a unique tempered correction.
If it is zero, choose the correction to be zero and retain the
discrepancy in the normal map.  Induction gives measurable tempered
coefficients at each finite order.  There are no resonances after
$b>d_{\max}/d_{\min}$, so the resulting normal map $N_v$ is a
finite weight-preserving polynomial.  Its linear diagonal is $A_v$,
and the weight induction just given supplies its polynomial inverse.

The construction respects the tower without choosing complements
inside an equal-weight block.  Project the degree-$b$ equation to
a lower level.  Since raw projection is linear and intertwines the
derivatives, this is precisely the lower-level homological equation.
On each nonresonant part uniqueness in
Lemma~\ref{nf:homological} identifies the projected correction with
the lower correction; on each resonant part both corrections are
zero and the projected discrepancy is retained.  Induction proves
compatibility of all formal changes and normal maps.

We now pass from a formal change to an exact smooth change.  This
step is included to specify the domains and finite-order estimates;
it is the measurable contraction mechanism underlying the usual
nonstationary normal-form results
\citep{Melnick2019,KalininSadovskaya2017}.
Let $\Psi_v^{[b]}$ be the polynomial truncation of the constructed
formal change through degree $b$.  Once $b$ exceeds the largest
resonant degree, the defect
\begin{equation}\label{nf:flat-defect}
 R_v^{[b]}=\Psi_{fv}^{[b]}G_v-N_v\Psi_v^{[b]}
\end{equation}
vanishes through degree $b$.  On a fixed raw domain its $C^q$ norms
are tempered for every fixed $b,q$.  This follows directly from the
raw bounds and the finitely many tempered polynomial coefficients.

Choose a positive radius $r_v$ on which
$\|DG_v\|\leq e^{-\kappa}$ for a fixed $\kappa>0$.
The raw second-derivative bounds and tempered norm comparison allow
$r_v$ to be chosen tempered.  Slowly decrease these radii so that
\[
 r_{fv}\geq e^{-\delta}r_v,\qquad 0<\delta<\kappa.
\]
Then $G_v(B(0,r_v))\subset B(0,r_{fv})$.  On this forward
admissible family of balls,
\begin{equation}\label{nf:raw-contraction}
 \|G_v^{(n)}(z)\|\leq e^{-\kappa n}\|z\|.
\end{equation}
For each fixed $q$, the derivatives through order $q$ of
$G_v^{(n)}$ have tempered bounds with arbitrarily small exponential
slack.  To see this directly, differentiate
$G_v^{(n+1)}=G_{f^nv}\circ G_v^{(n)}$.  At derivative order $q$ the
term containing $D^qG_v^{(n)}$ is multiplied by
$DG_{f^nv}$; all other terms are products of derivatives of lower
order and a tempered fixed-order derivative of $G_{f^nv}$.
Induction on $q$ and summation of the contracting scalar recurrence
give the claim.  One can obtain decay as well, but the stated
subexponential upper bounds suffice below.

There is a finite exponent $K$, depending on this state level but
independent of $b$, such that all coefficients of $N_v^{(-n)}$ are
bounded by $C_\epsilon(v)e^{(K+\epsilon)n}$.  Here $N_v^{(-n)}$
denotes the inverse of $N_v^{(n)}$ from the space at $f^nv$ back to
the space at $v$.  To check this bound, consider the finite-dimensional
space of coordinate monomials of a given weight $d$.  Substitution
by $N$ is triangular by ordinary degree; every diagonal tensor
block has the single exponent $-d$.  The triangular-array estimate
therefore bounds forward substitution by $e^{(-d+\epsilon)n}$ and
its inverse by $e^{(d+\epsilon)n}$, with tempered prefactors.
Only weights up to $d_{\max}$ occur in the coordinate functions of
$N^{\pm n}$, so one can take $K=d_{\max}$, with additional slack.
The same exponential bound controls any fixed derivatives of these
polynomials on a fixed bounded set, since their degrees are bounded
independently of $n$.

Define
\begin{equation}\label{nf:limit-formula}
 Q_{v,n}^{[b]}=N_v^{(-n)}\circ\Psi_{f^nv}^{[b]}
                          \circ G_v^{(n)}.
\end{equation}
Consecutive terms differ by applying $N_v^{(-n-1)}$ to two arguments
whose difference is $R_{f^nv}^{[b]}\circ G_v^{(n)}$.
Both arguments and each of their fixed derivatives are bounded
with arbitrarily small exponential slack; the arguments themselves
converge to zero.  Taylor's formula for
\eqref{nf:flat-defect}, \eqref{nf:raw-contraction}, and the chain rule
give, for $b\geq k$,
\begin{equation}\label{nf:defect-estimate}
 \bigl\|R_{f^nv}^{[b]}\circ G_v^{(n)}\bigr\|_{C^k}
 \leq C_{b,k,\epsilon}(v)
       e^{[-\kappa(b+1-k)+\epsilon]n}.
\end{equation}
Indeed, a term with $q\leq k$ derivatives falling on $R^{[b]}$
contains its Taylor factor of order $b+1-q$, evaluated at the
exponentially contracting argument, and only a fixed finite number
of derivatives of $G_v^{(n)}$.  All remaining costs are tempered.
Using the integral difference formula for the polynomial
$N_v^{(-n-1)}$ and differentiating it through order $k$ now yields
\begin{equation}\label{nf:telescoping}
 \|Q_{v,n+1}^{[b]}-Q_{v,n}^{[b]}\|_{C^k}
 \leq C_{b,k,\epsilon}(v)
 e^{[K-\kappa(b+1-k)+\epsilon]n}.
\end{equation}
The slack has been enlarged to absorb the finitely many chain-rule
factors; it is still arbitrarily small for fixed $b,k$.

Choose $b_0$ so that $\kappa(b_0+1)>K$.  The series of differences
converges in $C^0$ on the forward admissible domain.  For any fixed
$k$, choose a larger $b$ with $\kappa(b+1-k)>K$; it converges in
$C^k$ there.  These limits agree.  In fact
$\Psi^{[b']}-\Psi^{[b]}$ vanishes through degree $b$ when $b'>b$,
and the same estimate with $k=0$ shows
$Q_{v,n}^{[b']}-Q_{v,n}^{[b]}\to0$ whenever $b\geq b_0$.
Thus the common limit $\Psi_v$ is smooth.  Increasing $b$ does not
introduce a new convergence-domain assumption: each
$\Psi^{[b]}$ is a polynomial, and every fixed derivative of the raw
maps is bounded on the original raw domains.  On compact subsets
of the forward interior, the preceding estimates apply directly.
If a fixed-order auxiliary estimate uses a smaller tempered
radius at future centers, \eqref{nf:raw-contraction} and slow
variation eventually put the arguments inside it.

The identity relating $Q_{v,n+1}^{[b]}$ and
$Q_{fv,n}^{[b]}\circ G_v$ proves \eqref{nf:conjugacy} on taking
limits.  For any fixed Taylor order $a$, take $b\geq a$ sufficiently
large for $C^a$ convergence.  Every telescoping difference has
zero Taylor jet through degree $b$, so $j^a\Psi_v=j^a\Psi_v^{[b]}$.
In particular $D\Psi_v(0)=\id$.  Summing
\eqref{nf:telescoping} gives tempered fixed-order bounds.
The inverse function theorem, first with the $C^2$ bound to obtain
a tempered injectivity radius and then differentiated to any fixed
order, gives the stated inverse bounds.  For example, shrink until
$\|D\Psi_v-\id\|\leq1/2$; inverse derivatives are then finite sums
of products of derivatives of $\Psi_v$ and inverse matrices of norm
at most two.  These estimates use only finitely many tempered
quantities at every fixed order.  Finally, projection commutes with
each expression \eqref{nf:limit-formula}, after choosing a sufficiently
large common truncation for the two levels.  It therefore commutes
with their germ limits.
\end{proof}

\subsection{Fiber polynomiality and finite weight blocks}

\begin{lemma}[Analyticity only in the hidden variables]
\label{nf:fiber-analyticity}
At every fixed state level, the exact coordinate map $\Psi_{v,j}$
is polynomial in the raw hidden variables at fixed raw base $a$,
with smooth coefficients in $a$.  Write the normalized variables
as $(z_0,\zeta)$, where $z_0$ is the level-zero coordinate.  The
local inverse of $\Psi_{v,j}$ is smooth in $z_0$ and real analytic
in $\zeta$.  Every fixed derivative in $z_0$ and $\zeta$ has this
same real-analytic dependence on $\zeta$.  The assertion makes no
analyticity claim in the base or in the center $v$.
\end{lemma}

\begin{proof}
Fix a truncation $b_0$ for which the $C^0$ limit in
\eqref{nf:limit-formula} converges.  By
Lemma~\ref{nf:raw-dynamics}, at fixed raw $a$ the map
$G_v^{(n)}$ has polynomial hidden degree bounded by a number $D_j$
independent of $n$, and its base component is independent of the
hidden variables.  The polynomial $\Psi_{f^nv}^{[b_0]}$ has degree
at most $b_0$.  The inverse normal iterate is weight preserving,
so it has degree at most $L_j=d_{\max}/d_{\min}$, rounded down.
Consequently the hidden degree of every $Q_{v,n}^{[b_0]}$ is at
most $D_jb_0L_j$.  Polynomials of bounded degree form a closed
finite-dimensional space under uniform convergence on a ball:
evaluation at a fixed unisolvent finite collection of points
recovers all coefficients by an invertible matrix.  The $C^0$
limit is therefore polynomial of this bounded hidden degree.
Its coefficients are smooth in $a$, either by the same finite
evaluation formula on a smaller product neighborhood or by the
already established smoothness of $\Psi_v$.

Compatibility with level zero gives the triangular expression
\[
 \Psi_v(a,\eta)=(\psi_v(a),F_v(a,\eta)),
\]
where $\psi_v$ is a smooth local diffeomorphism and $F_v$ is
polynomial in $\eta$.  At fixed $a$, its hidden derivative is
invertible near the central state.  The real analytic inverse
function theorem in the hidden variables therefore gives an
analytic solution $\eta=H_v(a,\zeta)$ of
$F_v(a,\eta)=\zeta$.  The ordinary smooth parameter inverse
function theorem identifies this with a jointly smooth function.
Moreover, differentiating the implicit identity in a base variable
gives
\[
 \partial_aH_v
 =-(D_\eta F_v)^{-1}\,\partial_aF_v,
\]
evaluated at $(a,H_v(a,\zeta))$.  The right side is analytic in
$\zeta$ because $F_v$ and every base derivative of $F_v$ are
polynomial in $\eta$.  Repeated implicit differentiation expresses
each mixed derivative as a finite sum of products of such
polynomials, lower derivatives of $H_v$, and
$(D_\eta F_v)^{-1}$.  Induction proves the same analyticity for all
fixed derivatives.  Substituting $a=\psi_v^{-1}(z_0)$ proves the
statement in normalized base coordinates.  If desired, a common
small complex hidden neighborhood for each fixed real base point
is obtained by polynomial complexification and implicit inversion;
no complexification of the base is involved.
\end{proof}

\begin{corollary}[Finite arrays and higher kernels]
\label{nf:finite-weight}
In either normalized tower, the independent coordinates introduced
at level $j$ have weights at least $j\chi_*$ for a fixed
$\chi_*>0$.  A prescribed upper weight bound consequently uses
only finitely many independent tower coordinates and finitely many
monomials.  On the array of monomials of a fixed weight $d$,
normal substitution has the single exponent $-d$, and inverse
substitution has exponent $d$, with two-sided estimates and
arbitrarily small exponential slack.  These estimates continue to
hold for finite tensor arrays and their invariant subquotients,
with the corresponding sums and differences of weights.
\end{corollary}

\begin{proof}
The assertion about new coordinates is the quotient assertion in
Lemma~\ref{nf:raw-dynamics}, preserved by the compatible grading in
Lemma~\ref{nf:normalization}.  Positivity bounds the ordinary degree
of a monomial of weight at most $D$ by $D/\chi_*$, and only levels
$j\leq D/\chi_*$ can contribute independent variables.  At these
finitely many levels the dimensions are finite.  Substitution is
triangular by ordinary degree on each fixed-weight array; its
diagonal consists of tensor products of the derivative blocks
with total exponent $-d$.  The triangular-array estimates prove
the asserted forward and inverse bounds.  Tensor operations add
or subtract the indicated exponents.  Restriction preserves an
upper bound and the inverse upper bound supplies the lower bound;
the same two estimates pass to quotient norms.  This proves the
subquotient assertion.
\end{proof}

The corollary concerns finite arrays.  It will allow a later argument
to fix a weight cutoff, choose a sufficiently large spatial Taylor
order, and then pass to an orbit limit.  It does not identify an
infinite formal series with a convergent spatial germ.
